\section*{अध्याय \ref{ch:b3:probability} --- प्रायिकता: आधार
और बृहत् संख्याओं का नियम}

\begin{solution}{exo:b3:probability:1}
(क) $u \in \intoo01$ के लिए: $\P(F(X) \leq u) = \P(X \leq
F^{-1}(u)) = F(F^{-1}(u)) = u$ (\omterm{def:b3:topology:continuity}{सांतत्य} और कड़ी एकदिष्टता
$F$ को $\intoo01$ पर एकैकी आच्छादक बना देती हैं, जिससे
$\{F(X) \leq u\} = \{X \leq F^{-1}(u)\}$): अतः $F(X)$
एकसमान है। विपरीत दिशा में $\P(G(U) \leq t) = \P(U \leq F(t))
= F(t)$: अर्थात् किसी \omterm{def:b3:probability:space}{बंटन नियम} का अनुकरण करने के लिए किसी
एकसमान प्रतिदर्श पर प्रतिलोम बंटन फलन लगाइए।

(ख) $Y = X^2$, जहाँ $X$ $\intcc{-1}1$ पर एकसमान है: $t \in
\intcc01$ के लिए $F_Y(t) = \P(-\sqrt t \leq X \leq \sqrt t) =
\sqrt t$: अतः \omterm{ex:b3:lebesgue:gamma}{घनत्व} $\frac1{2\sqrt t}\mathbf 1_{\intoo01}$।
और $\P\bigl(-\frac1\lambda\ln U \leq t\bigr) = \P(U \geq
\eu^{-\lambda t}) = 1 - \eu^{-\lambda t}$: अर्थात् चरघातांकी
$\mathcal E(\lambda)$ --- प्रतिलोमन क्रिया में।
\end{solution}

\begin{solution}{exo:b3:probability:2}
(क) प्वासों: $\E X = \sum_{k\geq0}k\,\eu^{-\lambda}
\frac{\lambda^k}{k!} = \lambda$, $\E[X(X-1)] = \lambda^2$,
अतः $\V = \lambda^2 + \lambda - \lambda^2 = \lambda$।
ज्यामितीय ($\P(X = k) = p(1-p)^{k-1}$): $\E X = \frac1p$, $\V
= \frac{1-p}{p^2}$ (ज्यामितीय श्रेणी का दो बार अवकलन कीजिए)।

(ख) $G(t) = \P(T > t)$ अवर्धमान है और $G(0^+)\dots$
$G \colon \intco0\infty \to \intoc01$; और स्मृतिहीनता का पाठ
$G(t + s) = G(t)G(s)$ है। तब $G(n t) = G(t)^n$ तथा $G(t/n) =
G(t)^{1/n}$: अतः परिमेय $q \geq 0$ के लिए $G(q) = G(1)^q$;
$G(1) = \eu^{-\lambda}$ लिखिए ($\in \intoo01$: $G(1) = 1$
$G \equiv 1$ अनिवार्य कर देता, जो किसी परिमित \omterm{def:b3:probability:space}{यादृच्छिक चर} के
लिए असंभव है; और $G(1) = 0$ परिकल्पना से अपवर्जित है), और
किसी भी $t$ को परिमेयों के बीच दबाइए (एकदिष्टता): $G(t) =
\eu^{-\lambda t}$ --- अर्थात् चरघातांकी \omterm{def:b3:probability:space}{बंटन नियम}। विलोम एक
परिकलन मात्र है।
\end{solution}

\begin{solution}{exo:b3:probability:3}
(क) $\sigma(f_i(X_i)) = f_i(X_i)^{-1}(\mathcal B) \subseteq
X_i^{-1}(\mathcal B) = \sigma(X_i)$ ($f_i$ बोरेल), और \omterm{def:b3:probability:independence}{स्वतंत्र}
$\sigma$-बीजगणितों के उप-$\sigma$-बीजगणित \omterm{def:b3:probability:independence}{स्वतंत्र} होते हैं
(परिभाषक सर्वसमिका तब तो और भी सरलता से लागू होती है)।

(ख) $\sigma(A_i) = \{\varnothing, A_i, A_i^c, \Omega\} =
\sigma(A_i^c) = \sigma(\mathbf 1_{A_i})$: तीनों कथन एक ही
$\sigma$-बीजगणितों के \omterm{def:b3:probability:independence}{स्वातंत्र्य} का दावा करते हैं। ($A_i$ पर
का गुणनखंडन पूरकों तक फैलता है --- यह
\cref{def:b3:probability:independence} की तुल्यता के भीतर का
$\lambda$-निकाय तर्क है, अथवा सीधे समावेश-अपवर्जन।)

(ग) $\P(A) = \P(B) = \P(C) = \frac12$; और युग्मों पर $A\cap B
= A\cap C = B\cap C$: प्रत्येक प्रतिच्छेदन ``दोनों चित''
अथवा उसके अनुरूप है, जिसकी प्रायिकता $\frac14$ है: अतः
युग्मशः \omterm{def:b3:probability:independence}{स्वतंत्र}। पर $\P(A\cap B\cap C) =
\P(\text{चित-चित}) = \frac14 \neq
\frac18$: अतः \omterm{def:b3:probability:independence}{स्वतंत्र} नहीं --- $C$ $A$ और $B$ से निर्धारित
हो जाता है।
\end{solution}

\begin{solution}{exo:b3:probability:4}
(क) मान लीजिए पाठ $T$ की लंबाई $L$ है और $q = a^{-L}$ ($a$
वर्णमाला का आकार)। घटनाएँ $E_k = \{$स्थान $kL+1,
\dots, (k+1)L$ $T$ की वर्तनी करते हैं$\}$ \omterm{def:b3:probability:independence}{स्वतंत्र} हैं
(\omterm{def:b3:probability:independence}{स्वतंत्र} सम-बंटित अक्षरों के असंयुक्त खंड), और प्रत्येक की
प्रायिकता $q > 0$ है: अतः $\sum\P(E_k)
= \infty$, और बोरेल--कांतेली (2) से लगभग निश्चित रूप से अनंत
बार प्रकटन मिलते हैं।

(ख) ऊपरी: $\P\bigl(R_n \geq (1+\varepsilon)\log_2n\bigr)
\leq 2^{-(1+\varepsilon)\log_2n} = n^{-(1+\varepsilon)}$,
जिसका योग परिमित है: अतः बोरेल--कांतेली (1) से लगभग निश्चित
रूप से ऐसे $n$ परिमित ही हैं। निचली: असंयुक्त खंड भरिए ---
$j$-वाँ खंड लंबाई $\ell_j = \lceil\log_2s_j\rceil$ का, जो
$s_j = \sum_{i<j}\ell_i$ से आरंभ होता है; घटनाएँ ``खंड $j$
सर्वथा एक है'' \omterm{def:b3:probability:independence}{स्वतंत्र} हैं और उनकी प्रायिकता $2^{-\ell_j}
\asymp \frac1{s_j} \asymp \frac1{j\log_2 j}$ है, जिसका योग
अपसरित होता है: अतः बोरेल--कांतेली (2) से सर्वथा-एक वाले खंड
अनंत बार मिलते हैं, अर्थात् $R_{s_j} \geq \log_2 s_j$ अनंत
बार। दोनों मिलकर: प्रथम $n$ अंकों में सबसे लंबी लड़ी लगभग
निश्चित रूप से $(1 + o(1))\log_2n$ है।
\end{solution}

\begin{solution}{exo:b3:probability:5}
(क) $R = \bigl(\limsup\abs{X_n}^{1/n}\bigr)^{-1}$ परिमित
संख्या के $X_n$ बदलने पर अपरिवर्तित रहता है: अतः प्रत्येक $N$
के लिए $R$ $\sigma(X_N, X_{N+1}, \dots)$-मापनीय है, अर्थात्
पुच्छ-मापनीय। तब प्रत्येक घटना $\{R \leq c\}$ की प्रायिकता
$0$ या $1$ है (\cref{thm:b3:probability:zeroone}), इसलिए $R$
का बंटन फलन केवल मान $0, 1$ लेता है: वह किसी एक ही बिंदु $c_0
\in \intcc0{+\infty}$ पर उछलता है, और लगभग निश्चित रूप से
$R = c_0$।

(ख) $\sum X_n$ का तथा $\frac{S_n}n$ का अभिसरण परिमित संख्या
के पदों के बदलने से अप्रभावित रहता है (दूसरे के लिए: बदले हुए
पद $O(1/n) \to 0$ का योगदान करते हैं): अतः दोनों पुच्छ घटनाएँ
हैं; \omterm{thm:b3:probability:zeroone}{शून्य--एक नियम}।

(ग) $\{X_1 > 0\}$ $X_1$ पर निर्भर है: \omterm{def:b3:probability:independence}{स्वतंत्र} सम-बंटित
चिह्नों ($\P(X_1 = \pm1) = \frac12$) के लिए उसकी प्रायिकता
$\frac12 \notin \{0,1\}$ है --- इसमें कोई विरोधाभास नहीं, वह
पुच्छ घटना है ही नहीं।
\end{solution}

\begin{solution}{exo:b3:probability:6}
$(\intcc01, \lambda)$ पर काम कीजिए।
(क) टाइपराइटर $\mathbf 1_{I_n}$
(\cref{exo:b3:lp:3}): $\norm{X_n}_p^p = \lambda(I_n) \to 0$
(सभी $p < \infty$), अतः प्रायिकता में भी; पर प्रत्येक $\omega$
पर मान $0$ और $1$ दोनों बार-बार लौटते हैं: अतः कहीं भी
बिंदुशः अभिसरण नहीं।
(ख) $X_n = n\mathbf 1_{\intoo0{1/n}} \to 0$ $0$ के बाहर, पर
$\norm{X_n}_p \geq n^{1 - 1/p} \geq 1$।
(ग) $X_n = \sqrt n\,\mathbf 1_{\intoo0{1/n}}$: $\E\abs{X_n} =
n^{-1/2} \to 0$, $\E X_n^2 = 1$।
(घ) किसी भी उपानुक्रम में से (प्रायिकता में अभिसरण) कोई ऐसा
आगे का उपानुक्रम निकालिए जो लगभग निश्चित रूप से अभिसरित हो
(\cref{prop:b3:probability:modes}(ग)); प्रभावी अभिसरण उसके
अनुदिश $L^1$ अभिसरण दे देता है, और सीमा \emph{वही} $X$ रहती
है। अतः संख्यात्मक अनुक्रम $\E\abs{X_n - X}$ के प्रत्येक
उपानुक्रम का कोई उप-उपानुक्रम $0$ की ओर जाता है: इसलिए समूचा
अनुक्रम $0$ की ओर जाता है।
\end{solution}

\begin{solution}{exo:b3:probability:7}
(क) $\hat p_n = \frac{S_n}n$, जहाँ $S_n$ द्विपद है: $\V(\hat
p_n) = \frac{p(1-p)}n \leq \frac1{4n}$, और चेबिशेव
(\cref{prop:b3:probability:markov}) वह परिबंध दे देती है।
(ख) $\frac1{4n(0.03)^2} \leq 0.05$ हल कीजिए: $n \geq
\frac{1}{4\cdot0.0009\cdot0.05} \approx 5556$। उसी
प्रत्याभूति के लिए केंद्रीय सीमा प्रमेय $n \approx 1070$
उचित ठहरा देगी: चेबिशेव अपनी व्यापकता की क़ीमत लगभग $5$ के
गुणक से चुकाती है।
\end{solution}

\begin{solution}{exo:b3:probability:8}
(क) यदि $S_n \sim \mathcal B(n, x)$ हो, तो अंतरण प्रमेय
$\E\bigl[f(\frac{S_n}n)\bigr] =
\sum_k\binom nkx^k(1-x)^{n-k}f(\frac kn) = B_nf(x)$ देती है।

(ख)--(ग) मान लीजिए $\omega = \omega_f$ \omterm{def:b3:topology:continuity}{सांतत्य} मापांक है
($\abs{f(u) - f(v)} \leq \omega(\abs{u - v})$, और चरणों को
जोड़ने से $\omega(c
\delta) \leq (1 + c)\,\omega(\delta)$)। तब किसी भी $\delta >
0$ के लिए
\[
\abs{f(u) - f(x)} \leq \Bigl(1 + \frac{(u -
x)^2}{\delta^2}\Bigr)\omega(\delta)
\]
(यदि $\abs{u - x} \leq \delta$ हो तो स्पष्ट; अन्यथा
$\omega(\abs{u-x}) \leq (1 + \frac{\abs{u-x}}\delta)
\omega(\delta) \leq (1 + \frac{(u-x)^2}{\delta^2})
\omega(\delta)$)। $u = \frac{S_n}n$ पर \omterm{def:b3:probability:space}{प्रत्याशा} लीजिए:
\[
\abs{B_nf(x) - f(x)} \leq
\Bigl(1 + \frac{\V(S_n/n)}{\delta^2}\Bigr)\omega(\delta)
\leq \Bigl(1 + \frac{1}{4n\delta^2}\Bigr)\omega(\delta) ;
\]
और $\delta = n^{-1/2}$ के साथ: $\norm{B_nf - f}_\infty \leq
\frac54\,\omega\bigl(n^{-1/2}\bigr) \leq
\frac32\,\omega\bigl(n^{-1/2}\bigr) \to 0$ (\omterm{def:b3:topology:compact}{संहत} पर एकसमान
\omterm{def:b3:topology:continuity}{सांतत्य}): अर्थात् वाइरश्ट्रास प्रमेय की एक प्रायिकता-आधारित
उपपत्ति, स्पष्ट और एकसमान दर के साथ।
\end{solution}

\begin{solution}{exo:b3:probability:9}
(क) $k - 1$ प्रकार संग्रहीत हो चुकने के बाद प्रत्येक प्राचय
प्रायिकता $p_k = \frac{n-k+1}n$ के साथ नया होता है: अतः
$\tau_k$ ज्यामितीय $(p_k)$ है, और $\tau_k$ \omterm{def:b3:probability:independence}{स्वतंत्र} हैं
(प्राचय \omterm{def:b3:probability:independence}{स्वतंत्र} हैं)। योग लीजिए: $\E T_n =
\sum_k\frac n{n-k+1} = nH_n \sim n\ln n$;
$\V(T_n) = \sum\frac{1 - p_k}{p_k^2} \leq
n^2\sum_{j=1}^n\frac1{j^2} \leq \frac{\pi^2}6n^2$।

(ख) चेबिशेव: $\P\bigl(\abs{T_n - nH_n} \geq \varepsilon
n\ln n\bigr) \leq \frac{\pi^2n^2/6}{\varepsilon^2n^2\ln^2n}
\to 0$, और $\frac{nH_n}{n\ln n} \to 1$: अतः प्रायिकता में
$\frac{T_n}{n\ln n} \to 1$।

(ग) संघ परिबंध: $T_n > t$ का अर्थ है कि $\lceil t\rceil$
प्राचयों के बाद कोई प्रकार अनदेखा रह गया, अतः $\P(T_n > t)
\leq n(1 - \frac1n)^{t} \leq n\,\eu^{-t/n}$; और $t = \beta
n\ln n$ पर: $\leq n^{1 - \beta}$। $\beta > 1$ के लिए $\sum_m
2^{m(1-\beta)} < \infty$: अतः बोरेल--कांतेली $n = 2^m$ के
अनुदिश लगभग निश्चित रूप से अंततः $T_n \leq \beta n\ln n$ दे
देती है --- विशेषतः कथनानुसार प्रत्येक $\beta > 2$ के लिए
(उपानुक्रम के अनुदिश कोई भी $\beta > 1$ काम कर जाता है)।
\end{solution}

\begin{solution}{exo:b3:probability:10}
(क) सम-सूचकांक वाले अंक $(b_{2k})_k$ \omterm{def:b3:probability:independence}{स्वतंत्र} सम-बंटित
निष्पक्ष बिट हैं (\omterm{def:b3:probability:independence}{स्वतंत्र} अंक-कुल का एक उपकुल), अतः $U =
\sum_kb_{2k}2^{-k}$ प्रत्येक द्विआधारी अंतराल को उसकी सही
प्रायिकता देता है (जैसा
\cref{thm:b3:probability:existence} में): इसलिए एकसमान; और
वैसे ही $V$; और $(U, V)$ असंयुक्त अंक-खंडों पर निर्भर हैं:
अतः \omterm{def:b3:probability:independence}{स्वतंत्र} (द्विआधारी आयतों पर गुणनखंडन, फिर दिन्किन)।

(ख) $\Phi(\omega) = (U(\omega), V(\omega))$ \omterm{def:b3:lebesgue:measurable}{मापनीय} है और
$\Phi_*\lambda = \lambda\otimes\lambda = \lambda_2$
(द्विआधारी आयतों पर सहमति $+$ अद्वितीयता)। अंकों को गूँथने
से कोई प्रतिलोम मिलता है, जो किसी भी गुणनखंड के द्विआधारी
परिमेयों के (अकिंचन) समुच्चय के बाहर परिभाषित है: अर्थात्
$\intcc01$ और $\intcc01^2$ के पूर्ण-माप उपसमुच्चयों के बीच
कोई माप-संरक्षी एकैकी आच्छादन। पियानो
(\cref{pb:b3:topology:1}) से तुलना कीजिए: वहाँ \omterm{def:b3:topology:continuity}{सांतत्य} ने
एकैकीपन के बिना आच्छादकता अनिवार्य कर दी थी; और यहाँ \omterm{def:b3:topology:continuity}{सांतत्य}
छोड़कर केवल \omterm{def:b3:lebesgue:measurable}{मापनीयता} रखने पर कोई माप-तुल्याकारिता मिल जाती है
--- विमा \omterm{def:b3:measure:measure}{माप} सिद्धांत को नहीं दिखती, \omterm{def:b3:topology:topology}{संस्थिति} को दिखती है।
\end{solution}

\begin{solution}{exo:b3:probability:11}
(क) बंटन के \omterm{def:b3:topology:continuity}{संतत} होने से बराबरियाँ अकिंचन घटनाएँ हो जाती हैं
(जैसा अध्याय के क्रम-सांख्यिकी तर्कों में), और $(X_1, \dots,
X_n)$ के $n!$ सापेक्ष क्रम विनिमेय हैं, अतः समान रूप से
संभाव्य। $R_n = 1$ का अर्थ है कि उच्चिष्ठ अंतिम स्थान पर
बैठता है: जिसकी प्रायिकता $\frac{(n-1)!}{n!} = \frac1n$ है।

(ख) $n$ स्थिर कीजिए और $X_1, \dots, X_{n-1}$ के सापेक्ष
क्रम पर प्रतिबंधन कीजिए: $X_n$ को $n$ संभव कोटि-स्थानों में
डालना एकसमान है और उस क्रम से \omterm{def:b3:probability:independence}{स्वतंत्र} ($n$-बहुक की
विनिमेयता)। अतः $R_n$ (घटना ``$X_n$ शीर्ष स्थान लेता है'')
समूचे अभिलेख इतिहास $(R_1, \dots, R_{n-1})$ से \omterm{def:b3:probability:independence}{स्वतंत्र} है,
जो प्रथम $n - 1$ चरों के सापेक्ष क्रम का फलन है। आगमन से
पूर्ण स्वातंत्र्य मिल जाता है, जिसमें $\P(R_n = 1) =
\frac1n$।

(ग) \omterm{def:b3:probability:independence}{स्वातंत्र्य} के साथ $\sum\P(R_n = 1) = \sum\frac1n =
\infty$: अतः बोरेल--कांतेली का दूसरा अर्ध भाग लगभग निश्चित
रूप से अनंत बार अभिलेख दे देता है (अभिलेख कभी रुकते नहीं ---
पर वे लघुगणकीय रूप से विरल होते जाते हैं:
$\E[\#\text{अभिलेख} \leq n] = H_n \approx \ln n$)। क्रमागत
अभिलेख: $\P(R_n = R_{n+1}
= 1) = \frac1{n(n+1)}$ (\omterm{def:b3:probability:independence}{स्वातंत्र्य}), और
\[
\sum_n\frac1{n(n+1)} = \sum_n\Bigl(\frac1n -
\frac1{n+1}\Bigr) = 1 < \infty :
\]
अतः बोरेल--कांतेली का पहला अर्ध भाग लागू होता है --- लगभग
निश्चित रूप से क्रमागत-अभिलेख युग्म परिमित बार ही घटित होते
हैं।
\end{solution}

\begin{solution}{exo:b3:probability:12}
(क) स्थान $i \leq n$ से आरंभ होने वाली लंबाई $\ell$ की किसी
लड़ी की प्रायिकता $2^{-\ell}$ है; संघ परिबंध: $\P(L_n \geq
\ell) \leq n2^{-\ell}$। $\ell_n = (1 +
\varepsilon)\log_2n$ के साथ: $\P(L_n \geq \ell_n) \leq
n^{-\varepsilon}$। $n = 2^k$ के अनुदिश
$\sum_k2^{-k\varepsilon} < \infty$, अतः लगभग निश्चित रूप से
अंततः $L_{2^k} < (1+\varepsilon)k$ (बोरेल--कांतेली); और
व्यापक $n$ के लिए $2^{k-1} < n \leq 2^k$ चुनिए तथा $L_n$ की
एकदिष्टता और $\log_22^{k-1} \leq \log_2n$ का उपयोग कीजिए:
$L_n \leq L_{2^k}
< (1 + \varepsilon)k \leq (1 + \varepsilon)\frac{k}{k-1}
\log_2n$, और यह अतिरिक्त गुणक $\varepsilon$ को थोड़ा बढ़ाकर
आत्मसात कर लिया जाता है।

(ख) $\ell = \lceil(1 - \varepsilon)\log_2n\rceil$ और
$m = \lfloor n/\ell\rfloor$ असंयुक्त खंडों के साथ: खंड
\omterm{def:b3:probability:independence}{स्वतंत्र} हैं, प्रत्येक प्रायिकता $2^{-\ell}
\geq n^{-(1-\varepsilon)}/2$ के साथ सर्वथा चित है, अतः किसी
अचर $c > 0$ और बड़े $n$ के लिए
\[
\P(L_n < \ell) \leq \bigl(1 - 2^{-\ell}\bigr)^{m}
\leq \exp\bigl(-m2^{-\ell}\bigr)
\leq \exp\Bigl(-c\,\frac{n^{\varepsilon}}{\log_2n}\Bigr)
\]
होता है। ये प्रायिकताएँ $n = 2^k$ के अनुदिश योग्य हैं
(वस्तुतः सभी $n$ के अनुदिश): अतः बोरेल--कांतेली लगभग निश्चित
रूप से अंततः $L_n \geq (1 -
\varepsilon)\log_2n$ दे देती है ($2^k$ के बीच की जगह
एकदिष्टता (क) की भाँति निरापद रूप से भर देती है)।

(ग) दोनों परिबंध अनुक्रम $\varepsilon = \frac1j$ के अनुदिश
लीजिए और गणनीय संख्या की पूर्ण-माप घटनाओं का सर्वनिष्ठ
लीजिए: लगभग निश्चित रूप से $\frac{L_n}{\log_2n} \to 1$। $n =
10^6$ के लिए: $\log_2n
\approx 19.9$ --- अर्थात् लगभग $20$ चितों की लड़ी कोई
संदेहास्पद विसंगति नहीं, गणितीय निश्चितता है, और उसका न होना
इस बात का प्रमाण है कि किसी मनुष्य ने ``यादृच्छिकता'' गढ़ी
है (मनुष्य लगातार $5$ या $6$ से अधिक चित लिखने का साहस
विरले ही करते हैं)।
\end{solution}

\begin{solution}{pb:b3:probability:1}
\textbf{1.} $X_n^{\pm}$ $X_n$ के बोरेल फलन हैं: अतः वे
युग्मशः \omterm{def:b3:probability:independence}{स्वतंत्र}
(\cref{exo:b3:probability:3}(क)) और सम-बंटित तथा \omterm{def:b3:lebesgue:l1}{समाकलनीय}
बने रहते हैं, और $\E X_1 = \E X_1^+ - \E
X_1^-$। यदि प्रमेय अऋणात्मक चरों के लिए सत्य हो, तो उसे दोनों
अर्ध भागों पर लगाकर घटाइए:
लगभग निश्चित रूप से $\frac{S_n}n = \frac{S_n^+}n -
\frac{S_n^-}n \to \E X_1^+ - \E X_1^- = m$।

\textbf{2.} $\P(X_n \neq Y_n) = \P(X_n > n) = \P(X_1 > n)$
(सम \omterm{def:b3:probability:space}{बंटन नियम}), और $\sum_n\P(X_1 > n) \leq \sum_n\P(X_1
\geq n) \leq \E X_1 < \infty$
(\cref{exo:b3:product:3}(क))। बोरेल--कांतेली (1): लगभग
निश्चित रूप से सभी बड़े $n$ के लिए $X_n = Y_n$, अतः $S_n -
S_n^*$ $n$ में अंततः अचर है: इसलिए लगभग निश्चित रूप से
$\frac{S_n - S_n^*}n \to 0$, और दोनों मानकीकृत योगों का
अनंतस्पर्शी व्यवहार एक ही है।

\textbf{3.} $X_1\mathbf 1_{X_1 \leq n} \nearrow X_1$: अतः
एकदिष्ट अभिसरण प्रमेय $\E Y_n \to m$ देती है; और किसी
अभिसारी अनुक्रम के चेज़ारो माध्य उसी सीमा की ओर अभिसरित होते
हैं: $\frac{\E S_n^*}n =
\frac1n\sum_{k\leq n}\E Y_k \to m$। अतः केवल लगभग निश्चित
रूप से $\frac{S^*_n - \E S^*_n}{n} \to 0$ सिद्ध करना पर्याप्त
है।

\textbf{4.} $\V(Y_n) \leq \E Y_n^2 = \E[X_1^2\mathbf
1_{X_1\leq n}]$। श्रेणियों के लिए टोनेली से,
\[
\sum_n\frac{\E[X_1^2\mathbf 1_{X_1\leq n}]}{n^2}
= \E\Bigl[X_1^2\!\!\sum_{n \geq \max(X_1, 1)}\!\frac1{n^2}
\Bigr]
\leq \E\Bigl[X_1^2\cdot\frac{4}{\max(X_1,1)}\Bigr]
\leq 4\,\E[X_1] < \infty,
\]
जिसमें $x \geq 1$ के लिए $\sum_{n\geq x}n^{-2} \leq \frac4x$
का उपयोग हुआ ($x \geq 2$ के लिए: $\leq \frac1{x-1} \leq
\frac2x$; और $1
\leq x < 2$ के लिए: $\leq \frac{\pi^2}6 \leq \frac4x$,
क्योंकि $\frac4x > 2$), तथा दोनों स्थितियों $X_1 \gtrless 1$
में $X_1^2/\max(X_1, 1) \leq X_1$।

\textbf{5.} युग्मशः \omterm{def:b3:probability:independence}{स्वातंत्र्य} से $i \neq j$ के लिए
$\E[(Y_i - \E
Y_i)(Y_j - \E Y_j)] = 0$ (दो चरों वाला गुणनफल सूत्र), अतः
प्रसरण जुड़ जाते हैं: $\V(S^*_k) =
\sum_{n\leq k}\V(Y_n)$। प्रत्येक $k_j$ पर चेबिशेव लगाकर योग
लीजिए:
\[
\sum_j\P\Bigl(\abs{S^*_{k_j} - \E S^*_{k_j}} \geq
\varepsilon k_j\Bigr)
\leq \frac1{\varepsilon^2}\sum_j\frac1{k_j^2}\sum_{n\leq
k_j}\V(Y_n)
= \frac1{\varepsilon^2}\sum_n\V(Y_n)\!\!\sum_{j : k_j\geq
n}\!\frac1{k_j^2}
\]
(अऋणात्मक द्विक श्रेणी के लिए टोनेली)।

\textbf{6.} $k_j = \lfloor\alpha^j\rfloor \geq
\frac{\alpha^j}2$ (यह $\alpha^j \geq 1$ होते ही वैध है,
अर्थात् सभी $j \geq 0$ के लिए: $x \geq 1$ पर $\lfloor
x\rfloor \geq \frac x2$)। अतः
\[
\sum_{j : k_j \geq n}\frac1{k_j^2}
\leq 4\sum_{j : \alpha^j \geq n}\alpha^{-2j}
\leq \frac{4}{1 - \alpha^{-2}}\cdot\frac1{n^2}
= \frac{C_\alpha}{n^2},
\]
($\alpha^j \geq n$ वाले पहले $j$ से आरंभ होती ज्यामितीय
श्रेणी)। प्रश्न 4--5 के साथ मिलाने पर द्विक योग परिमित है;
और प्रत्येक परिमेय $\varepsilon$ के लिए लगाई गई तथा
सर्वनिष्ठ ली गई बोरेल--कांतेली (1) लगभग निश्चित रूप से
$\frac{S^*_{k_j} - \E S^*_{k_j}}{k_j} \to 0$ दे देती है, और
प्रश्न 3 के साथ: लगभग निश्चित रूप से
$\frac{S^*_{k_j}}{k_j} \to m$।

\textbf{7.} $Y_n \geq 0$ से $n \mapsto S^*_n$ अवघटमान हो
जाता है: अतः $k_j \leq n \leq k_{j+1}$ के लिए
\[
\frac{S^*_{k_j}}{k_{j+1}} \leq \frac{S^*_n}{n} \leq
\frac{S^*_{k_{j+1}}}{k_j},
\]
जो $\frac{k_j}{k_{j+1}}$ और $\frac{k_{j+1}}{k_j}$ अंतर्निविष्ट
करने के बाद कथित सैंडविच ही है। चूँकि
$\frac{k_{j+1}}{k_j} \to \alpha$, इसलिए प्रश्न 6 लगभग निश्चित
रूप से यह दे देता है:
\[
\frac m\alpha \leq \liminf_n\frac{S^*_n}n \leq
\limsup_n\frac{S^*_n}n \leq \alpha m .
\]

\textbf{8.} प्रश्न 7 को $\alpha = 1 + \frac1p$, $p
\in \N^*$ के लिए लगाइए: गणनीय संख्या की प्रायिकता-एक घटनाएँ;
उनके सर्वनिष्ठ पर $p \to \infty$ लेने से लगभग निश्चित रूप से
$\lim\frac{S^*_n}n = m$। प्रश्न 1--3 के साथ, लगभग निश्चित
रूप से $\frac{S_n}n \to \E X_1$: अर्थात् युग्मशः \omterm{def:b3:probability:independence}{स्वातंत्र्य}
के अंतर्गत बृहत् संख्याओं का प्रबल नियम।

\textbf{9.} स्वातंत्र्य-प्रकार की परिकल्पनाएँ तीन बार आईं:
(1) प्रसरणों की योज्यता (प्रश्न 5) --- यहाँ युग्मशः पर्याप्त
है; (2) सम बंटन, कर्तन वाले योगों में (प्रश्न 2) और माध्य
के परिकलन में (प्रश्न 3) --- यहाँ \omterm{def:b3:probability:independence}{स्वातंत्र्य} लगा ही नहीं;
(3) बोरेल--कांतेली (1) (प्रश्न 2 और 6) --- यह किसी भी
\omterm{def:b3:probability:independence}{स्वातंत्र्य} के बिना वैध है। पूर्ण पारस्परिक स्वातंत्र्य कभी
लगाया ही नहीं गया: यही एतेमादी का प्रेक्षण है।

\textbf{10.} कोई आधार $b$ और कोई अंक $r$ स्थिर कीजिए। किसी
एकसमान $\omega$ के आधार-$b$ अंक $(d_k)$ $\{0, \dots, b-1\}$
पर \omterm{def:b3:probability:independence}{स्वतंत्र} सम-बंटित एकसमान हैं (अंक-सदिश का प्रत्येक मान
लंबाई $b^{-m}$ का अंतराल घेरता है: यह
\cref{thm:b3:probability:existence} का तर्क शब्दशः है)।
\omterm{def:b3:probability:independence}{स्वतंत्र} सम-बंटित परिबद्ध चरों $\mathbf
1_{d_k = r}$ पर लगाया गया प्रबल नियम देता है: लगभग निश्चित
रूप से अंक $r$ की आवृत्ति $\frac1b$ की ओर जाती है। गणनीय
संख्या के युग्मों $(b, r)$ पर सर्वनिष्ठ लेने पर: लगभग
प्रत्येक संख्या \emph{प्रत्येक आधार में सरल रूप से सामान्य}
है। कोई स्पष्ट असामान्य संख्या: $x = 0.100100100\ldots_2$
(एकों की आवृत्ति $\frac13 \neq \frac12$)। यह विरोधाभास नम्र
कर देने वाला है: लगभग सभी संख्याएँ सामान्य हैं, फिर भी
$\sqrt2$, $\eu$ या $\pi$ के लिए सामान्यता आज भी असिद्ध है ---
\omterm{def:b3:measure:measure}{माप} सिद्धांत प्रस्तुत किए बिना गिन लेता है।

\textbf{11.} \cref{exo:b3:probability:10} को दोहराने पर एक ही
एकसमान चर से $\intcc01^d$ पर \omterm{def:b3:probability:independence}{स्वतंत्र} सम-बंटित एकसमान
\emph{सदिशों} $U_k$ का अनुक्रम मिल जाता है
(\cref{thm:b3:probability:existence} के प्रत्येक $U_n$ के
अंक-समुच्चय को $d$ उपकुलों में बाँटिए)।
$g \in L^1(\intcc01^d)$ के लिए चर $g(U_k)$ \omterm{def:b3:probability:independence}{स्वतंत्र} सम-बंटित
और \omterm{def:b3:lebesgue:l1}{समाकलनीय} हैं, जिनका माध्य $\int g\,\dd\lambda_d$ है
(अंतरण): अतः प्रबल नियम देता है
\[
\frac1n\sum_{k=1}^ng(U_k)
\xrightarrow[n\to\infty]{\text{लगभग निश्चित}}
\int_{\intcc01^d}g\,\dd\lambda_d :
\]
अर्थात् \omterm{ex:b3:probability:sllnapps}{मोंते कार्लो} समाकलन प्रत्येक विमा में लगभग निश्चित
रूप से अभिसरित होता है --- त्रुटि का आकार केंद्रीय सीमा
प्रमेय (\cref{ch:b3:clt}) का विषय है।

\textbf{12.} मान लीजिए $A_k = \{\abs{S_k} \geq \varepsilon\}
\cap \bigcap_{j<k}\{\abs{S_j} < \varepsilon\}$: ये $A_k$
असंयुक्त हैं और उनका सम्मिलन $A = \{\max_{k\leq
n}\abs{S_k} \geq \varepsilon\}$ है। तब
\[
\E S_n^2 \geq \sum_{k=1}^n\E\bigl[S_n^2\mathbf 1_{A_k}\bigr]
= \sum_{k=1}^n\E\Bigl[\bigl(S_k^2 + 2S_k(S_n - S_k) + (S_n
- S_k)^2\bigr)\mathbf 1_{A_k}\Bigr]
\geq \sum_{k=1}^n\E\bigl[S_k^2\mathbf 1_{A_k}\bigr],
\]
क्योंकि मिश्रित पद लुप्त हो जाता है: $S_k\mathbf 1_{A_k}$
गुट $(Z_1, \dots, Z_k)$ का बोरेल फलन है, जो $(Z_{k+1},
\dots, Z_n)$ के फलन $S_n - S_k$ से \omterm{def:b3:probability:independence}{स्वतंत्र} है
(\cref{thm:b3:probability:independence}), अतः
$\E[S_k\mathbf 1_{A_k}(S_n - S_k)] = \E[S_k\mathbf
1_{A_k}]\,\E[S_n - S_k] = 0$। $A_k$ पर $S_k^2 \geq
\varepsilon^2$, जिससे $\E S_n^2 \geq
\varepsilon^2\sum_k\P(A_k) = \varepsilon^2\P(A)$; और $\E
S_n^2 = \sum_{k\leq n}\V(Z_k)$ (प्रसरण जुड़ते हैं)।
निर्णायक चरण गुणनखंडन है: $S_k\mathbf 1_{A_k}$ समूचे पहले
खंड का एक \emph{अरैखिक} फलन है, और दूसरे खंड से उसका
\omterm{def:b3:probability:independence}{स्वातंत्र्य} गुट-स्वातंत्र्य है --- $Z_i$ का युग्मशः
\omterm{def:b3:probability:independence}{स्वातंत्र्य} केवल युग्मों का सहसंबंध मिटाता है और इसे उचित
नहीं ठहरा पाता।

\textbf{13.} $N$ स्थिर कीजिए और प्रश्न 12 को $Z_{N+1},
\dots, Z_{N+m}$ पर लगाइए:
\[
\P\Bigl(\max_{N < k \leq N+m}\abs{S_k - S_N} >
\varepsilon\Bigr) \leq
\frac1{\varepsilon^2}\sum_{j=N+1}^{N+m}\V(Z_j) \leq
\frac{r_N}{\varepsilon^2},
\qquad r_N = \sum_{j>N}\V(Z_j) .
\]
ये घटनाएँ $m$ के साथ बढ़ती हैं; नीचे से \omterm{def:b3:topology:continuity}{सांतत्य}
$\P(\sup_{k>N}\abs{S_k - S_N} > \varepsilon) \leq
r_N/\varepsilon^2$ दे देता है, और परिकल्पना से $r_N \to 0$।
अतः प्रत्येक $p \in \N^*$ के लिए $\P\bigl(\bigcap_N\{\sup_{k>N}
\abs{S_k - S_N} > \frac1p\}\bigr) \leq \inf_Np^2r_N = 0$:
अर्थात् लगभग निश्चित रूप से प्रत्येक $p$ के लिए कोई ऐसा $N$
है जिसके साथ $\sup_{k>N}\abs{S_k - S_N} \leq \frac1p$ ($p$
पर गणनीय संख्या की प्रायिकता-एक घटनाओं का सर्वनिष्ठ लीजिए),
ताकि सभी $k, l > N$ के लिए $\abs{S_k -
S_l} \leq \frac2p$: अर्थात् आंशिक योग लगभग निश्चित रूप से
कोशी हैं, इसलिए लगभग निश्चित रूप से अभिसारी।

\textbf{14.} चर $Z_n = x_n\varepsilon_n$ \omterm{def:b3:probability:independence}{स्वतंत्र} हैं
(\omterm{def:b3:probability:independence}{स्वतंत्र} चरों के बोरेल फलन,
\cref{exo:b3:probability:3}(क)), केंद्रित हैं, और $\V(Z_n) =
x_n^2$: अतः $\sum_nx_n^2 < \infty$ होने पर प्रश्न 13 लागू
होता है और लगभग निश्चित अभिसरण दे देता है। व्यापक स्थिति में,
प्रत्येक $N$ के लिए $\sum_nx_n\varepsilon_n$ का अभिसरण
$\varepsilon_1, \dots, \varepsilon_N$ के मानों से अप्रभावित
रहता है: अतः अभिसरण की घटना \omterm{def:b3:probability:independence}{स्वतंत्र} अनुक्रम
$(\varepsilon_n)$ के पुच्छ $\sigma$-बीजगणित में है, इसलिए
कोल्मोगोरोव का \omterm{thm:b3:probability:zeroone}{शून्य--एक नियम}
(\cref{thm:b3:probability:zeroone}) उसकी प्रायिकता को $0$ या
$1$ होने पर विवश कर देता है।

\textbf{15.} (क) स्तर $\theta\E Z$ पर विभाजन करके ऊपरी
टुकड़े पर कोशी--श्वार्ज़ लगाने से,
\[
\E Z = \E\bigl[Z\mathbf 1_{Z \leq \theta\E Z}\bigr] +
\E\bigl[Z\mathbf 1_{Z > \theta\E Z}\bigr]
\leq \theta\,\E Z + \sqrt{\E Z^2}\,
\sqrt{\P(Z > \theta\E Z)} ,
\]
अतः $(1 - \theta)\E Z \leq \sqrt{\E Z^2\,\P(Z > \theta\E
Z)}$; अब वर्ग कीजिए। (ख) $T_n^4 =
\sum_{i,j,k,l}x_ix_jx_kx_l\,
\E[\varepsilon_i\varepsilon_j\varepsilon_k\varepsilon_l]$
का प्रसार कीजिए: \omterm{def:b3:probability:space}{प्रत्याशा} तब $1$ है जब सूचकांक युग्मों में
बँट जाएँ (चारों बराबर, अथवा दो भिन्न युग्म, जो $3$ प्रकार से
सज सकते हैं) और अन्यथा $0$ (कोई अयुग्मित चिह्न शून्य माध्य
का होता है और \omterm{def:b3:probability:independence}{स्वातंत्र्य} से बाहर निकल आता है)। अतः
\[
\E T_n^4 = \sum_kx_k^4 + 3\sum_{i\neq j}x_i^2x_j^2 =
3s_n^4 - 2\sum_kx_k^4 \leq 3s_n^4 .
\]
(ग) $Z = T_n^2$, $\E Z = s_n^2$, $\theta = \frac14$ के साथ
पेली--ज़िगमुंड:
\[
\P\Bigl(\abs{T_n} > \frac{s_n}2\Bigr) = \P\Bigl(T_n^2 >
\frac{s_n^2}4\Bigr) \geq \Bigl(\frac34\Bigr)^2
\frac{s_n^4}{3s_n^4} = \frac3{16} .
\]
यदि श्रेणी धनात्मक प्रायिकता के साथ अभिसरित होती, तो वह लगभग
निश्चित रूप से अभिसरित होती (प्रश्न 14), अतः लगभग निश्चित
रूप से $\sup_n\abs{T_n} < \infty$, और कोई $M$
$\P(\sup_n\abs{T_n} > M) < \frac3{16}$ संतुष्ट करता; पर
$s_n > 2M$ होते ही $\P(\abs{T_n} > M) \geq \P(\abs{T_n} >
\frac{s_n}2) \geq \frac3{16}$: विरोधाभास। अतः अपसरण लगभग
निश्चित है, और प्रश्न 14 के साथ द्विभाजन \omterm{def:b3:complete:complete}{पूर्ण} हो जाता है।

\textbf{16.} यहाँ $x_n = n^{-s}$ है और $\sum_nn^{-2s} <
\infty$ ठीक तब जब $s > \frac12$: अतः प्रश्न 14--15 से
$\sum_n\frac{\varepsilon_n}{n^s}$ लगभग निश्चित रूप से
अभिसरित होती है तभी और केवल तभी जब $s > \frac12$ हो ($s \leq
\frac12$ के लिए लगभग निश्चित अपसरण)। $\frac12 < s \leq 1$
के लिए यह अभिसरण कभी निरपेक्ष नहीं होता। तुलना शिक्षाप्रद
है: पूर्णतः एकांतर चिह्न प्रत्येक $s > 0$ पर सामर्थ्य $n^{-s}$
के साथ निरसन कर देते हैं, जबकि प्रारूपिक यादृच्छिक चिह्न
केवल वर्गमूल-सामर्थ्य पर निरसन करते हैं --- प्रश्न 21 का
यादृच्छिक भ्रमण $\sqrt n$ की भाँति बढ़ता है, और आबेल योगीकरण
ठीक उसी वृद्धि को $s > \frac12$ पर
$\sum\varepsilon_nn^{-s}$ के अभिसरण में बदल देता है।

\textbf{17.} (क) $\cosh\lambda =
\sum_k\frac{\lambda^{2k}}{(2k)!}$ और $\eu^{\lambda^2/2} =
\sum_k\frac{\lambda^{2k}}{2^kk!}$; और $(2k)! \geq 2^kk!$
पदशः लागू होता है, क्योंकि $\frac{(2k)!}{k!} =
\prod_{i=1}^k(k + i) \geq \prod_{i=1}^k(2i) = 2^kk!$
(प्रत्येक गुणनखंड $i \leq k$ पर $k + i \geq 2i$ संतुष्ट करता
है), जिससे वस्तुतः $(2k)! \geq 2^k(k!)^2 \geq 2^kk!$।
(ख) ध्यान दीजिए कि $a \leq 0 \leq b$ ($Z$ केंद्रित है), और
$z \mapsto \eu^{\lambda z}$ की उत्तलता से, $z \in \intcc ab$
के लिए:
\[
\eu^{\lambda z} \leq \frac{b - z}{b - a}\,\eu^{\lambda a} +
\frac{z - a}{b - a}\,\eu^{\lambda b},
\qquad\text{अतः}\qquad
\E\,\eu^{\lambda Z} \leq \frac{b\,\eu^{\lambda a} -
a\,\eu^{\lambda b}}{b - a}
= (1 - p)\eu^{-pt} + p\,\eu^{(1-p)t} = \eu^{\varphi(t)}
\]
जहाँ $p = \frac{-a}{b-a} \in \intcc01$, $t = \lambda(b -
a)$, $\varphi(t) = -pt + \log(1 - p + p\eu^t)$। तब
$\varphi(0) = 0$, $\varphi'(t) = -p + \frac{p\eu^t}{1 - p +
p\eu^t}$ $0$ पर लुप्त हो जाता है, और $\rho = \frac{p\eu^t}{1
- p + p\eu^t} \in \intcc01$ के साथ $\varphi''(t) = \rho(1 -
\rho) \leq \frac14$: अतः कोटि $2$ पर टेलर
$\varphi(t) \leq \frac{t^2}8 = \frac{\lambda^2(b-a)^2}8$ दे
देता है।

\textbf{18.} $\lambda > 0$ के लिए धनात्मक चर
$\eu^{\lambda(S_n - \E S_n)}$ पर लगाई गई मार्कोव
(\cref{prop:b3:probability:markov}) और \omterm{def:b3:probability:independence}{स्वतंत्र} चरों का
गुणनफल सूत्र देते हैं
\[
\P(S_n - \E S_n \geq t) \leq \eu^{-\lambda
t}\prod_{i=1}^n\E\,\eu^{\lambda(X_i - \E X_i)}
\leq \exp\Bigl(-\lambda t +
\frac{\lambda^2}8\sum_i(b_i - a_i)^2\Bigr),
\]
जो प्रत्येक केंद्रित $X_i - \E X_i
\in \intcc{a_i - \E X_i}{b_i - \E X_i}$ (चौड़ाई वही) पर लगाए
गए प्रश्न 17(ख) से मिलता है। $D =
\sum_i(b_i - a_i)^2$ के साथ घातांक को $\lambda = \frac{4t}{D}$
पर न्यूनतम करने से $-\frac{2t^2}D$ मिलता है। अधःपुच्छ के लिए
वही परिणाम $(-X_i)$ पर लगाइए।

\textbf{19.} $t = n\varepsilon$ और $D = n(b - a)^2$ लीजिए:
\[
\P\Bigl(\Bigl|\frac{S_n}n - m\Bigr| \geq \varepsilon\Bigr)
\leq 2\exp\Bigl(\frac{-2n^2\varepsilon^2}{n(b-a)^2}\Bigr)
= 2\exp\Bigl(\frac{-2n\varepsilon^2}{(b-a)^2}\Bigr),
\]
जिसका $n$ में योग परिमित है (ज्यामितीय-प्रकार की श्रेणी):
अतः बोरेल--कांतेली
(\cref{thm:b3:probability:borelcantelli}) से लगभग निश्चित
रूप से अंततः $\abs{\frac{S_n}n - m} < \varepsilon$; और
$\varepsilon = \frac1p$ पर सर्वनिष्ठ लेने से लगभग निश्चित
रूप से $\frac{S_n}n \to m$। तुलना: एतेमादी केवल $X_1 \in
L^1$ और युग्मशः \omterm{def:b3:probability:independence}{स्वातंत्र्य} माँगते हैं, पर कोई दर नहीं देते;
हॉफडिंग परिबद्धता और पूर्ण स्वातंत्र्य माँगती है, और प्रत्येक
परिमित $n$ पर स्पष्ट चरघातांकी प्रत्याभूति दे देती है ---
दोनों प्रमेयें एक ही सीमा के बारे में भिन्न प्रश्नों का उत्तर
देती हैं।

\textbf{20.} $g(U_k)$ \omterm{def:b3:probability:independence}{स्वतंत्र} सम-बंटित हैं, $\intcc01$ में
मान लेते हैं और उनका माध्य $\int g\,\dd\lambda_d$ है
(अंतरण), अतः $b_i - a_i = 1$, $t = n\varepsilon$ के साथ
प्रश्न 18 द्विपार्श्व परिबंध देता है, और
$2\eu^{-2n\varepsilon^2} \leq \delta$ तभी होता है जब
$\eu^{2n\varepsilon^2} \geq \frac2\delta$, अर्थात्
$n \geq \frac{\log(2/\delta)}{2\varepsilon^2}$।
$\varepsilon = \delta = 10^{-2}$ के लिए:
\[
n \geq \frac{\log 200}{2\cdot10^{-4}} =
\frac{5.2983\ldots}{0.0002} \approx 26\,492 :
\]
अर्थात् लगभग $26\,500$ प्रतिदर्श $99\%$ विश्वास के साथ $1\%$
यथार्थता की प्रत्याभूति दे देते हैं --- प्रत्येक विमा $d$ में,
और $\intcc01$ में मान लेने वाले प्रत्येक \omterm{def:b3:lebesgue:measurable}{मापनीय} समाकल्य के
लिए। प्रश्न 11 का प्रबल नियम अभिसरण का वचन तो देता था पर
परिमित $n$ पर कोई प्रत्याभूति नहीं; और प्रति अक्ष $k$ बिंदुओं
वाला कोई नियतांकी जालक $k^d$ मूल्यांकन माँगता है, जो $d$ में
चरघातांकी है। संकेंद्रण ही \omterm{ex:b3:probability:sllnapps}{मोंते कार्लो} को आशा के बजाय
\emph{विधि} बनाता है।

\textbf{21.} \omterm{def:b3:probability:independence}{स्वातंत्र्य} और गुणनफल सूत्र:
प्रश्न 17(क) से $\E\,\eu^{\lambda S_n} =
(\E\,\eu^{\lambda\varepsilon_1})^n
= (\cosh\lambda)^n \leq \eu^{n\lambda^2/2}$।
$\eu^{\lambda S_n}$ पर मार्कोव:
\[
\P(S_n \geq x) \leq \eu^{-\lambda x + n\lambda^2/2}
= \eu^{-x^2/(2n)}
\qquad\text{अनुकूलतम } \lambda = \frac xn \text{ पर},
\]
और $-S_n$ के लिए सममित परिबंध (वही \omterm{def:b3:probability:space}{बंटन नियम}) $\abs{S_n}$
वाले अचर को दुगुना कर देता है।

\textbf{22.} $\eta > 0$ स्थिर कीजिए और $n \geq 2$ के लिए
$x_n = (1 + \eta)\sqrt{2n\log n}$ रखिए:
\[
\P(\abs{S_n} \geq x_n) \leq 2\exp\bigl(-(1 +
\eta)^2\log n\bigr) = \frac{2}{n^{(1+\eta)^2}},
\]
जिसका योग परिमित है, क्योंकि $(1 + \eta)^2 > 1$।
बोरेल--कांतेली: लगभग निश्चित रूप से सभी बड़े $n$ के लिए
$\abs{S_n} < (1 + \eta)\sqrt{2n\log n}$, अतः लगभग निश्चित
रूप से $\limsup_n\frac{\abs{S_n}}{\sqrt{2n\log n}} \leq 1 +
\eta$; और $\eta = \frac1p$, $p \in \N^*$ वाली प्रायिकता-एक
घटनाओं का सर्वनिष्ठ लेने से कथन सिद्ध हो जाता है। आकार $n$
के भ्रमण का प्रारूपिक आयाम $\sqrt n$ (उसका प्रसरण) है, और
उसके निकृष्टतम उत्क्रमण भी उस मापक्रम से अधिक से अधिक
$\sqrt{2\log n}$ गुना ही आगे जाते हैं।

\textbf{23.} $n_j = 2^j$ और $x = (1 +
\eta)\sqrt{2n_j\log\log n_j}$ ($j \geq 2$ के लिए परिभाषित)
के साथ प्रश्न 21 देता है
\[
\P\bigl(S_{n_j} \geq x\bigr) \leq \exp\bigl(-(1 +
\eta)^2\log\log n_j\bigr) = (j\log 2)^{-(1+\eta)^2},
\]
जिसका $j$ में योग परिमित है, क्योंकि $(1 + \eta)^2 > 1$:
अतः बोरेल--कांतेली और $\eta = \frac1p$ लगभग निश्चित रूप से
$\limsup_jS_{n_j}/\sqrt{2n_j
\log\log n_j} \leq 1$ दे देते हैं। पूरे ऊपरी अर्ध भाग के लिए
जो शेष है वह जाँच-बिंदुओं के बीच का सेतु है: दिखाना पड़ता है
कि $\max_{n_j\leq n\leq n_{j+1}}S_n$
$(1+\eta)\sqrt{2n_j\log\log n_j}$ को परिमित बार ही पार करता
है, जिसके लिए \emph{गाउसीय} पुच्छों वाली कोई उच्चिष्ठ असमिका
चाहिए (लेवी की परावर्तन असमिका अथवा ओत्तावियानी की असमिका,
जो यहाँ सिद्ध नहीं की गईं)। प्रश्न 12 मात्रात्मक रूप से बहुत
दुर्बल है: वह प्रायिकता को
\[
\frac{n_j}{(1+\eta)^2\,2n_j\log\log n_j}
= \frac{1}{2(1+\eta)^2\log(j\log2)},
\]
से परिबद्ध करता है, जो $0$ की ओर तो जाता है पर $j$ में
\emph{योग्य नहीं} है: अतः बोरेल--कांतेली निष्कर्ष नहीं निकाल
सकती। पुनरावृत्त लघुगणक के नियम का निचला अर्ध भाग दूसरी
बोरेल--कांतेली प्रमेयिका को \omterm{def:b3:probability:independence}{स्वतंत्र} वृद्धियों $S_{n_{j+1}}
- S_{n_j}$ पर लगाता है, और उसके लिए गाउसीय-प्रकार के पुच्छों
के सुमेलित निचले परिबंध चाहिए। दोनों परिष्कार ईमानदार तृतीय
वर्ष की प्रायिकता हैं, एक पाठ्यक्रम और आगे; और यह समस्या बिना
सहायता के जो देती है वह ज्यामितीय समयों के अनुदिश
पुनरावृत्त-लघुगणक का ठीक-ठीक मापक्रम है।

\textbf{24.} प्रत्येक $\hat p_i$ $\intcc01$ में मान लेने
वाले और माध्य $\P(A_i)$ वाले $n$ \omterm{def:b3:probability:independence}{स्वतंत्र} सम-बंटित सूचक चरों
का औसत है: अतः हॉफडिंग $\P(\abs{\hat p_i - \P(A_i)} >
\varepsilon) \leq 2\eu^{-2n\varepsilon^2}$ देती है। संघ
परिबंध इसे $N$ से गुणा कर देता है। $2N\eu^{-2n\varepsilon^2}
\leq \delta$ हल कीजिए: $n \geq
\frac{\ln(2N/\delta)}{2\varepsilon^2}$। संख्यात्मक रूप से:
$\ln\frac{2\cdot10^6}{0.05} =
\ln(4\cdot10^7) \approx 17.5$, अतः $n \geq
\frac{17.5}{2\cdot10^{-4}} \approx 87\,600$: अर्थात्
\emph{एक} प्रायिकता का $\pm1\%$ तक आकलन लगभग $18\,500$
प्रतिदर्श माँगता है ($\ln(2/\delta)/2\varepsilon^2$), और
\emph{दस लाख} प्रायिकताओं का आकलन केवल $\approx 4.7$ गुना
अधिक --- एकसमानता की क़ीमत $\ln N$ है, $N$ नहीं: यही वह
प्रेक्षण है जो आनुभविक जोखिम न्यूनीकरण को, और उसके साथ यंत्र
अधिगम को, सांख्यिकीय रूप से संभव बनाता है।

\textbf{25.} चर $X_n = \frac{\varepsilon_n}
{n^\alpha}$ \omterm{def:b3:probability:independence}{स्वतंत्र}, केंद्रित और परिबद्ध हैं, जिनके लिए
$\sum_n\V(X_n) = \sum_nn^{-2\alpha}$। यदि $\alpha >
\frac12$: तो प्रसरण श्रेणी अभिसरित होती है, और एक-श्रेणी
प्रमेय (भाग VI) $\sum X_n$ का लगभग निश्चित अभिसरण दे देती
है। यदि $\alpha \leq \frac12$: तो प्रसरण श्रेणी अपसरित होती
है, और विलोम अर्ध भाग (भाग VI का पेली--ज़िगमुंड तर्क, जो लागू
है क्योंकि योज्य पद $1$ से परिबद्ध हैं) लगभग निश्चित अपसरण
दे देता है। निरपेक्ष अभिसरण $\sum n^{-\alpha} < \infty$
माँगता है: अर्थात् $\alpha > 1$। $\intoc{\frac12}1$ पर
श्रेणी लगभग निश्चित रूप से अभिसरित होती है जबकि
$\sum\abs{X_n} = \infty$ निश्चित रूप से: चिह्न प्रायिकता एक
के साथ निरसन का षड्यंत्र रच लेते हैं --- अर्थात् निरसन से
अभिसरण, जो किसी भी निरपेक्ष कसौटी को अदृश्य है, और (\omterm{thm:b3:probability:zeroone}{शून्य--एक
नियम} से) फिर भी नियतांकी निर्णय के साथ।
\end{solution}
